% --- Begin LaTeX fragment: Deterministic Replay via Monoid Homomorphism + Induction ---
\documentclass[11pt]{article}
\usepackage{amsmath,amssymb,amsthm,mathtools}
\usepackage{hyperref}
\usepackage{cleveref}

\newtheorem{definition}{Definition}[section]
\newtheorem{lemma}[definition]{Lemma}
\newtheorem{theorem}[definition]{Theorem}
\newtheorem{corollary}[definition]{Corollary}
\newtheorem{assumption}[definition]{Assumption}
\newtheorem{remark}[definition]{Remark}

\begin{document}

\section*{Deterministic Replay: Monoid Homomorphism and Induction Proof}

\begin{definition}[State space]
Let $\mathbb{B}=\{0,1\}$. Let $\Sigma$ denote the system state space (finite or otherwise well-typed).
\end{definition}

\begin{definition}[Events, event monoid]
Let $\mathcal{T}$ be a finite set of \emph{transition labels} and let $I$ be the set of \emph{input tokens} (including recorded external responses, binding/router identifiers, timestamps, etc.). Define the event set
\[
\mathcal{E} := \mathcal{T}\times I \times \mathcal{M},
\]
where $\mathcal{M}$ is an environment metadata space. The free monoid $\mathcal{E}^*$ consists of all finite sequences (strings) of events with concatenation $\cdot$ and empty word $\varepsilon$.
\end{definition}

\begin{definition}[Endomorphism monoid]
Let $\mathrm{End}(\Sigma)$ denote the monoid of endomorphisms on $\Sigma$ under composition $\circ$, with identity $\mathrm{id}_\Sigma$.
\end{definition}

\begin{definition}[Per-event semantics]
For each transition label $\tau\in\mathcal{T}$ define a \emph{pure} transition function
\[
\phi_\tau : \Sigma \times I \times R \times \mathcal{M} \to \Sigma,
\]
where $R$ denotes a deterministic random-output domain (the PRNG output space). For an event $e=(\tau,i,\eta)\in\mathcal{E}$ and a PRNG sample $r\in R$ the per-event endomorphism is
\[
F(e;r) : \Sigma \to \Sigma,\qquad F(e;r)(\sigma) := \phi_\tau(\sigma,i,r,\eta).
\]
\end{definition}

\begin{definition}[Telemetry record]
A telemetry record is a tuple $R=(g, s, w, \Delta)$ where
\begin{itemize}
  \item $g$ is a GUID trace identifier,
  \item $s$ is a cryptographic seed (integer),
  \item $w = e_1 e_2 \dots e_n \in \mathcal{E}^*$ is the ordered event log,
  \item $\Delta$ is an optional set of snapshot state(s) (checkpoint(s)).
\end{itemize}
Associated to seed $s$ and a deterministic PRNG is a sequence of outputs $(r_1,r_2,\dots)$.
\end{definition}

\begin{assumption}[Completeness of recorded nondeterminism]
All sources of nondeterminism present in the original execution are either:
\begin{enumerate}
  \item explicitly recorded in $w$ or $\Delta$, or
  \item deterministically derivable from the seed $s$ via the PRNG.
\end{enumerate}
\end{assumption}

\begin{assumption}[Purity of replay semantics]
During replay, the binding implementations used to compute $\phi_\tau$ are pure functions of their explicit inputs: prior state, input token, PRNG output, and environment metadata. External side effects are replaced by recorded responses or deterministic stubs.
\end{assumption}

\begin{definition}[Extension of $F$ to sequences]
Fix a record $R=(g,s,w,\Delta)$ and PRNG outputs $(r_1,r_2,\dots,r_n)$. For event $e_j$ define $F(e_j):=F(e_j;r_j)\in\mathrm{End}(\Sigma)$. Extend $F$ multiplicatively to sequences by
\[
F(e_1 e_2 \cdots e_n) := F(e_n)\circ F(e_{n-1})\circ \cdots \circ F(e_1).
\]
By convention, $F(\varepsilon):=\mathrm{id}_\Sigma$.
\end{definition}

\begin{lemma}[Monoid homomorphism]
The map
\[
F : (\mathcal{E}^*,\cdot,\varepsilon) \longrightarrow (\mathrm{End}(\Sigma),\circ,\mathrm{id}_\Sigma)
\]
defined above is a monoid homomorphism: for all $u,v\in\mathcal{E}^*$,
\[
F(u\cdot v) = F(v)\circ F(u),
\qquad F(\varepsilon)=\mathrm{id}_\Sigma.
\]
\end{lemma}

\begin{proof}
Immediate from the definition. The empty word maps to identity by definition. Concatenation $u\cdot v$ corresponds to the sequential application of event endomorphisms; by definition of composition ordering we have $F(u\cdot v)=F(v)\circ F(u)$. Hence $F$ preserves the monoid structure.
\end{proof}

\begin{theorem}[Deterministic Replay --- formal statement]
Let $R=(g,s,w,\Delta)$ be a telemetry record satisfying the stated assumptions. Let $w=e_1 e_2 \cdots e_n$ and let $(r_1,\dots,r_n)$ be the PRNG outputs deterministically generated from seed $s$ according to the agreed indexing convention. Then there exists a deterministic replay procedure $\mathrm{Replay}(R)$ that reconstructs the unique state sequence
\[
\sigma_0 \xrightarrow{e_1} \sigma_1 \xrightarrow{e_2} \cdots \xrightarrow{e_n} \sigma_n,
\]
where $\sigma_j = F(e_1\cdots e_j)(\sigma_0)$ and $\sigma_0$ equals the snapshot in $\Delta$ (if provided) or the canonical initial state.
\end{theorem}

\begin{proof}
We prove by induction on $n=|w|$ that replay reconstructs the same states as the original execution.

\paragraph{Base case ($n=0$).}
If $w=\varepsilon$ then $F(w)=\mathrm{id}_\Sigma$ and the replayed state sequence is the singleton $\{\sigma_0\}$ where $\sigma_0$ is the snapshot in $\Delta$ (or the canonical initial state). Equality with the original trivially holds.

\paragraph{Inductive hypothesis.}
Assume for some $k\ge 0$ that for any telemetry record whose event log length is $k$ the replay procedure reconstructs the original states $\sigma_0,\sigma_1,\dots,\sigma_k$ satisfying
\[
\sigma_j = F(e_1\cdots e_j)(\sigma_0)\quad\text{for }0\le j\le k.
\]

\paragraph{Inductive step ($k\to k+1$).}
Consider a record with $w=e_1\cdots e_k e_{k+1}$ and associated PRNG outputs $(r_1,\dots,r_{k+1})$. By the inductive hypothesis, replaying the prefix $e_1\cdots e_k$ produces the state $\sigma_k = F(e_1\cdots e_k)(\sigma_0)$. For step $k+1$ the replay computes
\[
\sigma_{k+1}' := F(e_{k+1};r_{k+1})(\sigma_k),
\]
where $F(e_{k+1};r_{k+1})$ is the per-event endomorphism computed with the recorded input token, environment metadata, and the PRNG output $r_{k+1}$. By the assumptions:
\begin{itemize}
  \item the PRNG output $r_{k+1}$ equals the original run's $r_{k+1}$ (determinism of PRNG given seed $s$),
  \item the input token and environment metadata match the original values (they are recorded),
  \item the per-event function $\phi_{\tau_{k+1}}$ is pure and deterministic for the given arguments.
\end{itemize}
Thus $\sigma_{k+1}'=\sigma_{k+1}$, where $\sigma_{k+1}$ is the state the original execution produced after event $e_{k+1}$. This completes the inductive step.

Therefore, by mathematical induction, replay reconstructs the same state sequence for all $n$.
\end{proof}

\begin{corollary}[Closure under concatenation]
If $R_1=(g_1,s_1,w_1,\Delta_1)$ and $R_2=(g_2,s_2,w_2,\Delta_2)$ are replayable traces and the concatenation $w_1\cdot w_2$ is provided with a consistent PRNG indexing scheme (or equivalent combined seed), then the concatenated trace is replayable and
\[
F(w_1\cdot w_2) = F(w_2)\circ F(w_1).
\]
\end{corollary}

\begin{remark}[Practical caveats]
The theorem is only valid under the stated assumptions. In practice, ensure:
\begin{itemize}
  \item explicit recording of external responses and wall-clock values if they influence $\phi_\tau$;
  \item deterministic handling (or recording) of concurrency interleavings (vector clocks, global counters, or total-order logs);
  \item PRNG outputs are indexed and/or recorded so the mapping event index $\mapsto r_j$ is unambiguous.
\end{itemize}
Failure to observe these requirements yields partial or no reproducibility.
\end{remark}

\end{document}
% --- End LaTeX fragment ---
